% ============================================================
% SECCIÓN 8: Despliegue y Operacionalización del Sistema
% ============================================================
\chapter{Despliegue y Operación}
\label{sec:despliegue}

% ─────────────────────────────────────────────────────────────
\section{Visión General del Despliegue}

El sistema Demeter se despliega en \textbf{cuatro entornos diferenciados}, cada uno
adaptado a un propósito distinto del ciclo de vida del proyecto. Los cuatro comparten
las mismas imágenes Docker (compiladas una sola vez por la pipeline de
\texttt{ghcr.io}); lo que cambia entre entornos son los \textit{overrides} de
configuración, los volúmenes y la forma de exposición al exterior.

\begin{table}[H]
    \centering
    \caption{Entornos de despliegue del sistema Demeter}
    \begin{tabularx}{\textwidth}{lXll}
        \toprule
        \textbf{Entorno} & \textbf{Propósito}                                                   & \textbf{Exposición}              & \textbf{Comando} \\
        \midrule
        Local (dev)      & Desarrollo con hot-reload, base de datos volátil.                    & \texttt{localhost:5173 / :8000}  & \texttt{make dev} \\
        Staging          & Pre-producción accesible vía dominio real, sin abrir puertos.        & Cloudflare Tunnel a \texttt{patata.monters.org} & \texttt{make staging} \\
        Producción (GCP) & Servicio público estable con Traefik como proxy inverso.            & \texttt{patata.monters.org}      & \texttt{make prod} (manual SSH) \\
        Raspberry Pi     & Pasarela de borde junto a los nodos ESP32; opcionalmente UI local.  & Red local + WSS al backend       & \texttt{make rpi-pull} \\
        \bottomrule
    \end{tabularx}
    \label{tab:envs}
\end{table}

\begin{designnote}
\textbf{Una imagen, varios entornos.} Las imágenes Docker se construyen una sola vez
en CI (push a \texttt{main}) y se publican en \texttt{ghcr.io}. Cada entorno hace
\texttt{pull} de la misma imagen y aplica un \texttt{docker-compose override} que
cambia variables de entorno, puertos y políticas de \texttt{restart}. Esto garantiza
que \textit{lo probado en staging es bit a bit lo desplegado en producción}.
\end{designnote}

La arquitectura desplegada en GCP se ilustra en la
Sección~\ref{fig:anx_backend_docker} del Anexo, y el equivalente en Raspberry Pi en
la Sección~\ref{fig:anx_rpi_docker}.

% ─────────────────────────────────────────────────────────────
\section{Despliegue del Firmware ESP32}
\label{sec:despliegue_firmware}

El firmware se construye y se sube a los nodos ESP32 mediante \textbf{PlatformIO}. La
configuración de los entornos vive en \texttt{Firmware/platformio.ini} y se ejecuta
desde la propia Raspberry Pi (que actúa como \textit{flashing host} con sus tres
puertos USB conectados a los ESP32 vía CDC).

\subsection{Catálogo de entornos}

\begin{table}[H]
    \centering
    \caption{Entornos de PlatformIO relevantes para v1.0}
    \begin{tabularx}{\textwidth}{lllX}
        \toprule
        \textbf{Env}                & \textbf{Rol}              & \textbf{Puerto USB}     & \textbf{Notas} \\
        \midrule
        \texttt{gateway}            & Gateway (Node 1)          & \texttt{/dev/ttyACM0}    & Receptor ESP-NOW + UART al host. \\
        \texttt{sensor\_cluster}    & Sensor cluster (Node 2)   & \texttt{/dev/ttyACM3}    & Sensor con N plantas, ESP-NOW puro. \\
        \texttt{actuador}           & Actuador (Node 3)         & \texttt{/dev/ttyACM2}    & Activación de relé de bomba. \\
        \texttt{sensor\_deepsleep}  & Sensor con Deep Sleep     & \texttt{/dev/ttyACM2}    & Variante de bajo consumo (validada en banco). \\
        \texttt{native}             & Tests unitarios PC        & ---                      & Banco de tests Unity (Sección~\ref{sec:testing_unit}). \\
        \texttt{get\_mac}           & Utilidad                  & \texttt{/dev/ttyACM1}    & Imprime la MAC del chip por serie. \\
        \bottomrule
    \end{tabularx}
    \label{tab:pio_envs}
\end{table}

El \texttt{platformio.ini} declara cuatro entornos adicionales (\texttt{sensor},
\texttt{sensor\_node\_3}, \texttt{sensor\_test}, \texttt{integration}) usados durante
el desarrollo y los tests; no forman parte del despliegue v1.0.

\subsection{Procedimiento de flasheo}

El flasheo es \textbf{USB local sobre la Raspberry Pi}. La identificación física de
los puertos se realiza con \texttt{pio device list} y se contrasta con las MACs
documentadas en \texttt{docs/flash\_instructions.md}.

\begin{lstlisting}[style=trama, caption={Flasheo de los tres nodos desde la Raspberry Pi}]
# 1. Identificar puertos disponibles
pio device list

# 2. Flashear los tres nodos (atajos del Makefile)
make fw-flash-gateway     # /dev/ttyACM0
make fw-flash-sensor      # /dev/ttyACM3
make fw-flash-actuador    # /dev/ttyACM2

# 3. Equivalente sin Makefile
cd Firmware
pio run -e gateway        -t upload --upload-port /dev/ttyACM0
pio run -e sensor_cluster -t upload --upload-port /dev/ttyACM3
pio run -e actuador       -t upload --upload-port /dev/ttyACM2
\end{lstlisting}

\subsection{Verificación post-flasheo}

Tras el flasheo, cada nodo debe emitir una traza identificable por el monitor serie:

\begin{lstlisting}[style=trama, caption={Verificación serie del gateway}]
make fw-monitor-gateway
# Salida esperada:
#   === DEMETER GATEWAY V2 (Direct Control) ===
#   [Setup] Ready.
#   >> [CLUSTER] Node 2: 4 entries
#       Plant 1: 22.50 C, 45% soil
\end{lstlisting}

\begin{warningbox}
\textbf{OTA no implementado en v1.0.} Las actualizaciones de firmware requieren
acceso físico a los puertos USB de la Raspberry Pi. Esta es una limitación
deliberada: implementar \textit{Over-The-Air} sobre ESP-NOW exige reservar
particiones, gestionar autenticación y manejar \textit{rollback} en caso de fallo,
trabajo que no se justifica para un sistema doméstico de un solo despliegue. Si en
una versión futura se distribuyera a múltiples invernaderos, OTA se convierte en un
requisito.
\end{warningbox}

% ─────────────────────────────────────────────────────────────
\section{Despliegue de la Raspberry Pi}
\label{sec:despliegue_rpi}

La Raspberry Pi cumple dos funciones simultáneas: \textbf{pasarela de borde}
(\textit{edge gateway}) entre los ESP32 y el backend, y opcionalmente
\textbf{servidor de UI local} para usuarios sin acceso a internet pero con pantalla
conectada al dispositivo.

\subsection{Provisioning del sistema}

Sobre una Raspberry Pi OS recién instalada, los pasos para preparar el host son:

\begin{lstlisting}[style=trama, caption={Provisioning de la Raspberry Pi}]
# 1. Instalar Docker y plugin compose
curl -fsSL https://get.docker.com | sh
sudo usermod -aG docker $USER && newgrp docker

# 2. Liberar el UART GPIO: deshabilitar consola serie del kernel
sudo raspi-config nonint do_serial 1   # console=disabled, hw=enabled
# Equivalente manual:
# - Editar /boot/firmware/cmdline.txt y eliminar 'console=serial0,115200'
# - Asegurar 'enable_uart=1' en /boot/firmware/config.txt
sudo reboot

# 3. Verificar que /dev/serial0 existe y el usuario pertenece a 'dialout'
ls -l /dev/serial0
groups | grep dialout || sudo usermod -aG dialout $USER

# 4. Clonar repositorio
git clone https://github.com/<usuario>/Proyecto_Demeter.git
cd Proyecto_Demeter
cp .env.rpi.example .env.rpi   # ajustar DEMETER_BACKEND_URL si fuera distinto
\end{lstlisting}

\subsection{Arranque de los servicios}

El despliegue en RPi se gestiona con \texttt{docker-compose.rpi.yml}, que define dos
servicios: \texttt{gateway} (FastAPI + bridge UART$\leftrightarrow$WS) y
\texttt{local-ui} (Nginx sirviendo la SPA React). Ambos corren en \texttt{network\_mode:
host} para evitar la sobrecarga del NAT de Docker en hardware ARM modesto.

\begin{lstlisting}[style=trama, caption={Arranque del stack en la Raspberry Pi}]
# Imagen pre-compilada (multi-arch arm64) desde GHCR (recomendado)
make rpi-pull

# O construir localmente si se han hecho cambios sin pasar por CI
make rpi-build

# Operación cotidiana
make rpi-logs        # logs en vivo
make rpi-status      # estado de contenedores y env vars activas
make rpi-shell       # entrar al contenedor 'demeter-gateway'
\end{lstlisting}

\subsection{Variables de entorno relevantes}

\begin{table}[H]
    \centering
    \caption{Variables de entorno consumidas por el contenedor \texttt{gateway}}
    \begin{tabularx}{\textwidth}{lXl}
        \toprule
        \textbf{Variable}              & \textbf{Función}                                       & \textbf{Default} \\
        \midrule
        \texttt{DEMETER\_BACKEND\_URL} & Host del backend (sin esquema).                        & \texttt{patata.monters.org} \\
        \texttt{DEMETER\_BACKEND\_PORT}& Puerto WSS hacia el backend.                          & \texttt{443} \\
        \texttt{DEMETER\_PORT}         & Path del dispositivo serie dentro del contenedor.     & \texttt{/dev/ttyS0} \\
        \texttt{DEMETER\_UART\_BAUD}   & Velocidad UART (debe coincidir con firmware Gateway). & \texttt{115200} \\
        \texttt{DEMETER\_LOG\_LEVEL}   & Verbosidad (\texttt{DEBUG}, \texttt{INFO}, \dots).    & \texttt{INFO} \\
        \texttt{HOST\_SERIAL\_PORT}    & Mapeo del dispositivo del host (\texttt{/dev/serial0}). & \texttt{/dev/serial0} \\
        \bottomrule
    \end{tabularx}
    \label{tab:env_rpi}
\end{table}

\subsection{Healthchecks y reinicio}

A diferencia del backend (Sección~\ref{sec:despliegue_backend}), los servicios de la
Raspberry Pi \textbf{sí incluyen healthchecks} en
\texttt{docker-compose.rpi.yml}: el \texttt{gateway} expone
\texttt{http://localhost:8001/api/health} y el \texttt{local-ui} responde en el puerto
80. La política de restart es \texttt{unless-stopped}, suficiente para reanudar tras
cortes de luz frecuentes en entornos no controlados.

% ─────────────────────────────────────────────────────────────
\section{Despliegue del Backend en GCP}
\label{sec:despliegue_backend}

El backend se despliega como un \textbf{stack Docker Compose} sobre una máquina
virtual de Google Cloud Platform. La VM no se aprovisiona automáticamente: se
asume creada manualmente con Ubuntu LTS y Docker instalado, y con el dominio
\texttt{patata.monters.org} apuntando a su IP pública.

\subsection{Composición del stack}

El stack en producción combina dos archivos Compose:

\begin{itemize}[leftmargin=*, label=--]
    \item \texttt{docker-compose.yml} (base): define los seis servicios
          (\texttt{frontend}, \texttt{api}, \texttt{math-worker}, \texttt{sequencer},
          \texttt{db}, \texttt{redis}) descritos en el Capítulo~\ref{sec:backend}.
    \item \texttt{docker-compose.prod.yml} (override): aplica \texttt{restart:
          always} a todos los servicios y añade \texttt{labels} de Traefik para
          enrutar \texttt{patata.monters.org} y sus paths \texttt{/api/} y
          \texttt{/ws/}.
\end{itemize}

\subsection{Procedimiento de despliegue}

El target \texttt{prod} del Makefile no orquesta el despliegue automáticamente:
imprime el procedimiento manual a seguir por SSH. Esta decisión es deliberada para
evitar despliegues accidentales contra producción desde un portátil de desarrollo.

\begin{lstlisting}[style=trama, caption={Procedimiento de despliegue en producción}]
# 1. Conexión al servidor
ssh danielmf31@patata.monters.org

# 2. En el servidor: actualizar código y configuración
cd ~/Proyecto_Demeter
git pull origin main
cp .env.prod .env

# 3. Levantar stack con override de prod
docker compose -f docker-compose.yml -f docker-compose.prod.yml up -d --build

# 4. Migraciones (idempotente)
docker exec demeter-api sh -c "cd /app/Software/Servidor/Backend && alembic upgrade head"

# 5. Seed inicial (solo en el primer despliegue)
docker exec demeter-api python -m BD.seed_data
\end{lstlisting}

\subsection{Reverse proxy y TLS}

\textbf{Traefik} actúa como proxy inverso único frente al stack. Termina TLS para
\texttt{patata.monters.org} (certificado gestionado por Let's Encrypt vía Traefik
ACME) y enruta:

\begin{itemize}[leftmargin=*, label=--]
    \item \texttt{patata.monters.org/} $\rightarrow$ \texttt{demeter-frontend:80}
    \item \texttt{patata.monters.org/api/...} $\rightarrow$ \texttt{demeter-api:8000}
    \item \texttt{patata.monters.org/ws/...} $\rightarrow$ \texttt{demeter-api:8000} (WebSocket)
\end{itemize}

Traefik no se incluye en \texttt{docker-compose.prod.yml}: se asume desplegado a
nivel de la VM, escuchando en \texttt{:80} y \texttt{:443} y conectado a la misma red
Docker que el stack Demeter mediante una red externa compartida.

\subsection{Variante staging con Cloudflare Tunnel}

Para validar cambios contra el dominio público sin abrir puertos en el portátil de
desarrollo, \texttt{docker-compose.staging.yml} añade un servicio
\texttt{cloudflared} que crea un túnel Cloudflare Zero Trust:

\begin{lstlisting}[style=trama, caption={Variante staging con túnel Cloudflare}]
# .env.staging debe definir CLOUDFLARE_TUNNEL_TOKEN obtenido en
# dash.cloudflare.com -> Zero Trust -> Tunnels.
make staging

# Acceso:
#   https://patata.monters.org      (vía túnel)
#   http://localhost:5175           (acceso directo sin túnel)
#   http://localhost:8001/api/docs  (Swagger directo)
\end{lstlisting}

Este modo es útil para demostraciones desde el portátil sin tocar la VM de
producción y para sesiones de QA en las que se quiere ejercitar el sistema con un
certificado TLS real sin comprometer datos productivos.

\subsection{Migraciones}

Las migraciones de PostgreSQL se gestionan con \textbf{Alembic}. El target
\texttt{make migrate} ejecuta \texttt{alembic upgrade head} dentro del contenedor
\texttt{demeter-api}. Las migraciones \textbf{inversas (downgrade) no están soportadas
oficialmente}: ante un fallo de migración se recomienda restaurar desde backup en
lugar de revertir esquema.

% ─────────────────────────────────────────────────────────────
\section{Pipeline de Integración y Entrega Continua (CI/CD)}
\label{sec:despliegue_cicd}

El proyecto cuenta con dos workflows de GitHub Actions ya operativos en
\texttt{.github/workflows/}, uno para validación continua y otro para construcción y
publicación de imágenes Docker.

\subsection{Workflow de validación: \texttt{ci\_develop.yml}}

Disparado en \texttt{push} y \texttt{pull\_request} contra la rama \texttt{develop}.
Su único job, \texttt{test-all}, prepara un entorno con Python 3.12 y Node 20,
instala las dependencias de las cuatro suites del proyecto y ejecuta el target
\texttt{make test-all} (Capítulo~\ref{sec:testing}). Aplica caché para \texttt{pip}
y \texttt{\textasciitilde/.platformio} para acelerar las ejecuciones sucesivas.

\subsection{Workflow de entrega: \texttt{deploy\_main.yml}}

Disparado en \texttt{push} a la rama \texttt{main}. Ejecuta primero
\texttt{test-all} como gate y, si pasa, lanza el job \texttt{build-and-push} con una
estrategia de matriz que construye \textbf{cuatro imágenes Docker multi-arquitectura}:

\begin{table}[H]
    \centering
    \caption{Imágenes Docker publicadas por la pipeline}
    \begin{tabularx}{\textwidth}{lllX}
        \toprule
        \textbf{Imagen}    & \textbf{Plataforma} & \textbf{Dockerfile}                                          & \textbf{Build args} \\
        \midrule
        \texttt{frontend}  & \texttt{linux/amd64} & \texttt{Software/Servidor/docker/frontend/Dockerfile}        & \texttt{VITE\_API\_URL}, \texttt{VITE\_WS\_URL}, \texttt{VITE\_DEV\_LOGIN} \\
        \texttt{backend}   & \texttt{linux/amd64} & \texttt{Software/Servidor/docker/backend/Dockerfile}         & --- \\
        \texttt{gateway}   & \texttt{linux/arm64} & \texttt{Software/Raspberry/Dockerfile}                       & --- \\
        \texttt{local-ui}  & \texttt{linux/arm64} & \texttt{Software/Servidor/docker/frontend/Dockerfile.rpi}    & \texttt{VITE\_API\_URL=/api}, \texttt{VITE\_WS\_URL=/ws} \\
        \bottomrule
    \end{tabularx}
    \label{tab:ghcr_images}
\end{table}

Cada imagen se etiqueta con cinco patrones simultáneos: nombre de rama, hash corto
de commit, \texttt{latest}, versión semver (si el commit lleva tag) y la etiqueta
fija \texttt{0.4} (versión de la pipeline). Las imágenes ARM64 se construyen sobre
runners x86 mediante \texttt{QEMU}, evitando depender de runners self-hosted en
Raspberry Pi.

\begin{warningbox}
\textbf{Lo que el CI no hace.} La pipeline construye imágenes pero
\textbf{no despliega automáticamente}: el \texttt{pull} desde \texttt{ghcr.io} y el
arranque del stack se realiza manualmente desde la VM o la Raspberry Pi. La razón
es operativa: en un sistema con cinco macetas reales, una promoción automática a
producción durante un experimento en curso puede invalidar los datos del
investigador. Esta es una decisión deliberada, no una deuda técnica.

Tampoco se construye firmware en CI ni se ejecuta análisis de secretos
(\texttt{gitleaks}); ambas tareas están planificadas para versiones futuras.
\end{warningbox}

% ─────────────────────────────────────────────────────────────
\section{Operaciones Post-Despliegue}
\label{sec:despliegue_ops}

Una vez el sistema está corriendo, las operaciones recurrentes (logs, backups,
diagnóstico) se realizan vía Makefile o vía \texttt{docker compose} directamente.

\subsection{Logs}

\begin{lstlisting}[style=trama, caption={Atajos de logs por servicio}]
make logs           # Todos los servicios del stack activo (tail 50)
make logs-api       # Solo API FastAPI
make logs-db        # Solo PostgreSQL
make logs-redis     # Solo Redis
make logs-frontend  # Solo Nginx (frontend)
make rpi-logs       # Stack Raspberry Pi (gateway + local-ui)
\end{lstlisting}

\textbf{No hay centralización de logs} (Loki, ELK, Cloud Logging). Para un sistema
de cinco macetas y un solo administrador, \texttt{docker logs} con
\texttt{logrotate} sobre el daemon es suficiente. Si en el futuro se desplegaran
múltiples invernaderos, se introduciría una capa de agregación.

\subsection{Backups de PostgreSQL}

\textbf{No hay backups automatizados} en v1.0 (deuda declarada del proyecto). El
procedimiento manual recomendado es:

\begin{lstlisting}[style=trama, caption={Backup manual semanal de la base de datos}]
# Backup
docker exec demeter-db pg_dump -U postgres demeter_db \
    | gzip > backup-$(date +%Y%m%d).sql.gz

# Restore (en una VM nueva o tras desastre)
gunzip -c backup-YYYYMMDD.sql.gz \
    | docker exec -i demeter-db psql -U postgres demeter_db
\end{lstlisting}

Para v1.0 se planifica una entrada \texttt{cron} semanal que escriba a un USB
externo conectado a la Raspberry Pi (Semana 4 del plan de release).

\subsection{Diagnóstico y troubleshooting}

Los problemas operativos recurrentes del proyecto están documentados en archivos
dedicados dentro de \texttt{docs/}:

\begin{table}[H]
    \centering
    \caption{Documentos de troubleshooting por capa}
    \begin{tabularx}{\textwidth}{lX}
        \toprule
        \textbf{Documento}                         & \textbf{Cubre} \\
        \midrule
        \texttt{Backend\_Problemas\_y\_Soluciones.md}    & Errores comunes del API y la DB. \\
        \texttt{Firmware\_Problemas\_y\_Soluciones.md}   & Errores de flasheo, ESP-NOW y handshake. \\
        \texttt{Raspberry\_Problemas\_y\_Soluciones.md}  & UART, permisos serie, conexión WSS. \\
        \texttt{Error\_Redis\_Connection.md}             & Conexión de la API a Redis. \\
        \texttt{Error\_Frontend\_WS\_UART.md}            & Errores WS y UART vistos desde el frontend. \\
        \texttt{Guia\_Monitorizacion\_TUI.md}            & Monitorización ad hoc de PostgreSQL/Redis desde CLI. \\
        \bottomrule
    \end{tabularx}
    \label{tab:troubleshooting_docs}
\end{table}

% ─────────────────────────────────────────────────────────────
\section{Actualización y Rollback}

\subsection{Actualización}

\textbf{Staging:} \texttt{make staging-down \&\& git pull \&\& make staging}
(arranca tirando de las imágenes \texttt{:latest} de \texttt{ghcr.io} ya
construidas por CI).

\textbf{Producción:} ssh a la VM, \texttt{git pull origin main}, repetir el comando
\texttt{docker compose ... up -d --build}. Se asume ventana de mantenimiento
acordada; el sistema no es \textit{zero-downtime}.

\textbf{Raspberry Pi:} \texttt{make rpi-pull} hace \texttt{pull} de la imagen
ARM64 más reciente y reinicia el stack. Tiempo típico de corte: 5--15 segundos.

\subsection{Rollback}

Las imágenes en \texttt{ghcr.io} están etiquetadas con el hash de commit, el nombre
de la rama, el tag semver y la etiqueta fija \texttt{0.4}. Para volver a una
versión anterior basta con apuntar el override al tag deseado:

\begin{lstlisting}[style=trama, caption={Rollback a una versión anterior}]
# Identificar el tag al que volver (e.g. sha-abc1234 o un semver previo)
DEMETER_BACKEND_IMAGE=ghcr.io/danielmf31/proyecto_demeter/backend:sha-abc1234 \
DEMETER_FRONTEND_IMAGE=ghcr.io/danielmf31/proyecto_demeter/frontend:sha-abc1234 \
make staging
\end{lstlisting}

\begin{warningbox}
\textbf{Migraciones no reversibles.} Si entre la versión actual y la versión a la
que se quiere volver hay una migración Alembic destructiva (e.g.\ una columna
borrada), el rollback de la imagen no basta: hay que restaurar también la base de
datos desde el backup más reciente. Por esa razón se recomienda hacer
\texttt{pg\_dump} antes de cualquier despliegue de producción que incluya
migraciones.
\end{warningbox}

% ─────────────────────────────────────────────────────────────
\section{Resumen}

El despliegue v1.0 del sistema Demeter sostiene cuatro entornos
(\textit{dev}, \textit{staging}, \textit{producción}, \textit{Raspberry Pi}) sobre
una única familia de imágenes construidas por CI y publicadas en
\texttt{ghcr.io}. La operación se gobierna desde un \textbf{Makefile con
$\sim$50 targets} que cubren desde el flasheo del firmware hasta la promoción
manual a producción. Las decisiones de scope explícitas
(sin OTA, sin auto-deploy, sin observabilidad centralizada, sin backups
automáticos) corresponden a la realidad de un sistema doméstico con un
administrador único; cada una está documentada como tal y pasa a deuda formal si en
una futura versión se ampliara el alcance del despliegue.
